\section{结语}

综合全文讨论，可以看到，当前毕设真正面临的问题并不是系统是否已经具备足够多的功能，而是原有选题主线仍然停留在“生命科学多 Agent 助手”这一应用层叙事之中。这样的叙事虽然能够体现一定的系统复杂度，却很难保证研究边界的独立性，也难以在答辩中形成一个集中、清晰且具有算法性的核心问题。特别是在已有另一生命科学助手类毕设存在的情况下，如果继续沿原有路线扩展功能，即使系统加入更多门控、HITL 与恢复机制，也仍然容易被视为父集系统的增强版本，而非一个独立成立的计算机专业毕业设计。

基于这一判断，本文将课题目标重新定义为：面向高代价科学工作流的运行时治理控制问题，并以蛋白设计场景作为核心验证对象。在这一新的问题框架下，系统关注的重点不再是“如何充当一个更完整的生命科学助手”，而是“如何在部分可观测、代价不对称、可失败且可恢复的工作流中，依据运行时证据决定下一步如何推进、恢复、人工介入或终止”。这样的转换不仅提升了选题层级，也使课题更自然地进入控制、调度、约束决策和系统治理等计算机问题领域。

在多个候选方向中，\texttt{controller} 路线之所以被进一步收敛为 \texttt{Belief-State Controller}，核心原因在于它最适合将“难度、风险、证据充分性、恢复空间和人工介入价值”统一纳入一个动态更新的控制框架。相较于静态门控或启发式打分，\texttt{Belief-State Controller} 更强调系统对隐状态的持续估计，以及在此基础上的治理动作选择；相较于完整的 POMDP 或强化学习求解，它又能够在当前时间约束下通过可解释的降阶版本实现落地。因此，这一路线兼顾了理论依据、算法主线、系统兼容性与阶段性可实现性，是当前条件下最有希望收敛为完整毕业设计成果的方向。

进一步来看，本文也明确区分了“已有方法背景”和“本课题可能的创新空间”。Belief State、uncertainty-aware planning、cost-aware routing、stepwise value guidance 和多 Agent failure analysis 等思想在近两年研究中都已有明确积累，因此本课题的贡献不应被简单表述为“使用了 belief state”。真正可能成立的创新点，应落在新的问题形式化、结构化隐状态设计、基于运行时异构事件的 belief 更新，以及工作流级治理动作空间的组织方式上。与此同时，本课题的难点也并不在于代码规模，而在于如何从隐变量、异构观测、长期代价与协作噪声中提炼出一个足够清晰、足够可解释的控制问题，并在现有系统中验证这一控制层的真实价值。

从可行性角度看，本文最终认为该方向在当前时间节点下仍然具备现实落地条件。其原因不在于这一方向“足够简单”，而在于现有系统已经提供了足够重的 runtime 平台和实验基础，使得后续工作可以集中于控制层的研究问题，而不必再重复底层系统建设。只要研究边界被严格控制在“高代价科学工作流治理控制”这一层，并坚持以蛋白设计工作流作为单一核心验证场景，那么形成一个可解释的 \texttt{Belief-State Controller} 原型及其对照验证结果，仍然是现实且可管理的目标。

因此，本文的最终结论是：当前毕设最合理的收敛方向，不是继续扩展一个更大的生命科学助手系统，而是基于现有平台，围绕高代价科学工作流提出并验证一个 \texttt{Belief-State Controller}。如果这一方向能够顺利完成，其最终成果将不只是现有系统中的局部增强，而是一个边界清晰、主线明确、兼具算法性与系统落地价值的毕业设计课题。
